\section{The published concentration coordinates and an exact conic heat transport}
\label{sec:ns-conic-transport}

This section uses the coordinate formulas of the public manuscript
\emph{Finite Time Blowup for Navier--Stokes}, whose author line is
OpenAI. The source reading receipt records the preserved revisions and
the audited source pin. The manuscript's Theorem~1.1 claims finite-time
unbounded velocity for every viscosity $\nu>0$, zero initial velocity,
a smooth compactly supported force, fixed compact spatial support before
time one, and uniformly bounded kinetic energy. That global claim is
not independently verified in this section. All statements proved below
are explicit coordinate, differential-algebra, and scalar-germ identities.
In particular, the leading field appearing in the source is retained as
a leading field, rather than substituted for its completed construction.
The current 166-page PDF has SHA-256
\begin{center}\small\texttt{0e779481c4da40bd28d1e642e1d8ca57447d129610df28dfa5a11e9af8ae228f}.\end{center}
The earlier 165-page PDF has SHA-256
\begin{center}\small\texttt{8c8a94ad9ac824c8b605b9827cadf7beaca48bd10b380de3cfc872a2c37afa81}.\end{center}
The source audit compares the cited passages in both versions: the
coordinates are (3.2), current p.~7 and (4.1), current p.~24;
the differential identities are Lemma~4.1 and (4.2), current p.~25;
the profiles and balances are (4.3)--(4.7), current pp.~25--26;
the growth and viscosity formulas are (10.21)--(10.23), current
p.~124. The corresponding earlier passages are on pp.~7, 24--25,
and 124. The selected equations agree; complete revision equivalence
is not asserted. The preserved source commit is
\begin{center}\small\texttt{8937a8f4cbc7abaab5e9e97d1cc7f5d2319d9538}.\end{center}
The public manuscript supplies these source objects
\cite{OpenAINSBridge2026}. Its revised historical account credits
Diego C\'ordoba, Luis Mart\'inez-Zoroa, and the indicated joint work
with Fan Zheng for preceding forced-flow constructions.
Tristan Buckmaster's primary statement separately identifies his
joint work with Levent Alp\"oge on announced smooth-forcing
incompressible-porous-media, Boussinesq, and three-dimensional Euler
blowup, and describes an unfinished verification of a hypodissipative
Navier--Stokes claim \cite{BuckmasterStatementBridge2026}.
These are the stated attributions and mathematical scopes. Neither
source observation establishes priority or a technical dependency
of the positive-viscosity manuscript on that joint work. The new
conic calculations below are proved here from the displayed objects.

\subsection{The exact source coordinates and differential operators}
To prevent collisions of source notation, write $\rho$ for the source's
cylindrical radius $r$, $Q$ for its concentration variable $q$, $\xi$
for its similarity variable $\eta$, and $\sigma=\rho^2/2$ for its
squared-radius variable $s$. The arithmetic coordinates $r,s,b,q$ retain
their previous meanings. These are names for the original variables,
not changes of their values. Retain
\begin{gather*}
 A=\tfrac12+h,\qquad D=\tfrac12-h,\qquad \tau=1-t,\\
 z=Q^D\xi,\qquad \tau=Q(1-\xi^2),\qquad X=\frac{\rho^2}{2Q},\\
 d=1-\xi^2,\qquad \ell=1-2h\xi^2.
\end{gather*}
The coordinate identities hold for $0<h<1/2$, $\tau>0$,
$Q>0$, $-1<\xi<1$; the source's profile construction uses
$0<h<1/100$. Powers of $Q$ denote their positive real values.
For each $(z,\tau)\in\mathbb R\times(0,\infty)$ there is exactly one
$Q>|z|^{1/D}$ with
\[
 Q-z^2Q^{2h}=\tau.
\]
Indeed the left side is zero at the stated endpoint, tends to infinity,
and has derivative $1-2hz^2Q^{2h-1}\geq1-2h>0$ on this interval.
The implicit-function theorem and $\ell>0$ prove smoothness. Inverting
gives $\xi=zQ^{-D}$ and $\rho=\sqrt{2QX}$, the latter for $X>0$.
The time--axial Jacobian is
\[
 \det\frac{\partial(\tau,z)}{\partial(Q,\xi)}
 =Q^D(d+2D\xi^2)=Q^D\ell>0.
\]
Differentiating these exact relations at fixed physical coordinates gives
\begin{align*}
 Q_t&=-\ell^{-1},& \xi_t&=D\xi/(Q\ell),& X_t&=X/(Q\ell),\\
 Q_z&=2\xi Q^{1-D}/\ell,&
 \xi_z&=d/(Q^D\ell),& X_z&=-2\xi X/(Q^D\ell),\\
 Q_\rho&=0,&\xi_\rho&=0,&X_\rho&=\rho/Q.
\end{align*}
For example $Q_\tau=1/\ell$, $\xi_\tau=-D\xi/(Q\ell)$,
and $X_\tau=-X/(Q\ell)$, so the first line retains the time reversal.
The $\xi_z$ calculation uses $\ell-2D\xi^2=d$.

For any real $w$ and smooth profile $f(X,\xi)$ define
$\mathcal P_w f=Q^w f(X,\xi)$ and
\begin{align*}
 T_wf&=\ell^{-1}(-wf+D\xi f_\xi+Xf_X),\\
 Z_wf&=\ell^{-1}(2w\xi f+d f_\xi-2\xi Xf_X),\qquad
 Rf=2Xf_{XX}+2f_X.
\end{align*}
The chain and product rules, with the preceding six derivatives, give
\begin{equation}\label{eq:ns-profile-operators-exact}
 \partial_t\mathcal P_w=\mathcal P_{w-1}T_w,
 \qquad \partial_z\mathcal P_w=\mathcal P_{w-D}Z_w,
 \qquad
 (\partial_\rho^2+\rho^{-1}\partial_\rho)\mathcal P_w
       =\mathcal P_{w-1}R.
\end{equation}
For the last identity, the two terms from differentiating $\rho/Q$
and the cylindrical term each contribute $Q^{w-1}f_X$, and
$\rho^2Q^{w-2}f_{XX}=2XQ^{w-1}f_{XX}$. Thus all radial constants
are retained. Applying the axial identity twice gives
\begin{align}
 (\partial_t-\nu\Delta)\mathcal P_w f
 &=\mathcal P_{w-1}(T_w-\nu R)f
     -\nu\mathcal P_{w-1+2h}Z_{w-D}Z_wf,
       \label{eq:ns-profile-heat-full}\\
 -\tfrac14\partial_t^2\mathcal P_w f
 &=-\tfrac14\mathcal P_{w-2}T_{w-1}T_wf.
       \label{eq:ns-profile-arithmetic-full}
\end{align}
Here $\Delta$ is the three-dimensional scalar Laplacian on
axisymmetric functions. These equations keep both distinct similarity
weights $w-1$ and $w-1+2h$, as well as the arithmetic factor $-1/4$.

The actual source leading profiles satisfy
\[
 u^{(0)}_\theta=Q^{-A}E,\quad u^{(0)}_z=Q^{-A}U,
 \quad \rho u^{(0)}_\rho=V_0,\quad p^{(0)}=Q^{-2A}\Pi,
 \qquad E=C_0^{-1}\sqrt{2X}\,\phi,\quad C_0>1,
\]
where $C_0$ denotes the source's amplitude constant, not the GCT
coefficient. Its incompressibility and pressure identities are
\[
 V_0=\frac X\ell\left(2\xi U-2D\xi\mathcal A_XU
                   -d\partial_\xi\mathcal A_XU\right),\quad
 \mathcal A_XU=X^{-1}\int_0^XU(y,\xi)\,dy,\qquad
 \Pi_X=E^2/(2X).
\]
The first follows by integrating
$(V_0)_X=-Z_{-A}U$ from zero, using $A+D=1$ and
integrating $XU_X$ exactly. The second is the source radial-pressure
balance $\partial_\rho p^{(0)}=(u^{(0)}_\theta)^2/\rho$.
The regular axis interpretation is the one stated by the source:
$E/\sqrt{2X}$, $U$, and $V_0/X$ are smooth at $X=0$.

For this particular leading field its physical material derivative on
axisymmetric scalars has the exact representation
\begin{equation}\label{eq:ns-source-material-profile}
 \bigl(\partial_t+u^{(0)}_\rho\partial_\rho
                     +u^{(0)}_z\partial_z\bigr)\mathcal P_w f
   =\mathcal P_{w-1}M_wf,\qquad
 M_w=T_w+V_0\partial_X+U Z_w.
\end{equation}
Indeed $u^{(0)}_\rho\partial_\rho\mathcal P_wf
=Q^{w-1}V_0f_X$, while the axial term has weight
$w-A-D=w-1$. This proves the correspondence for the actual source
ansatz, including its nonlinear dependence on the chosen profiles.

For completeness this also determines its full momentum residual,
including the cylindrical frame terms. Put $f_\rho=V_0/\sqrt{2X}$
and $R_1=R-(2X)^{-1}$. On $X>0$ the three components of
$(\partial_t+u^{(0)}\cdot\nabla)u^{(0)}-\nu\Delta u^{(0)}
+\nabla p^{(0)}$ are
\begin{align*}
 \mathfrak r_\rho={}&\mathcal P_{-3/2}
       (M_{-1/2}f_\rho-\nu R_1f_\rho)
       -\nu\mathcal P_{-3/2+2h}Z_{-1/2-D}Z_{-1/2}f_\rho,\\
 \mathfrak r_\theta={}&\mathcal P_{-A-1}
       (M_{-A}E+V_0E/(2X)-\nu R_1E)
       -\nu\mathcal P_{-A-1+2h}Z_{-A-D}Z_{-A}E,\\
 \mathfrak r_z={}&\mathcal P_{-A-1}
       (M_{-A}U-\nu RU+Z_{-2A}\Pi)
       -\nu\mathcal P_{-A-1+2h}Z_{-A-D}Z_{-A}U.
\end{align*}
To verify these formulas, differentiation of the cylindrical unit
vectors gives the radial convective term $-(u_\theta^{(0)})^2/\rho$,
the azimuthal convective term $u_\rho^{(0)}u_\theta^{(0)}/\rho$,
and the radial and azimuthal vector-Laplacian terms
$-u_\rho^{(0)}/\rho^2$ and $-u_\theta^{(0)}/\rho^2$.
The radial convective term cancels $\partial_\rho p^{(0)}$ by
the exact displayed pressure balance. The azimuthal additional term
is $\mathcal P_{-A-1}(V_0E/(2X))$.
Finally $-2A-D=-A-1$ gives the axial pressure weight. Substituting
\eqref{eq:ns-profile-operators-exact} proves every remaining term.
No assertion that this residual vanishes is needed or made.

The source constructs its leading profiles at viscosity one. Its
actual passage to arbitrary $\nu>0$ is the spatial map
\[
 y=x/\sqrt\nu,\qquad
 u_\nu(x,t)=\sqrt\nu\,u(y,t),\quad
 p_\nu(x,t)=\nu p(y,t),\quad
 f_\nu(x,t)=\sqrt\nu\,f(y,t).
\]
Every momentum term acquires the factor $\sqrt\nu$: time
differentiation preserves the velocity prefactor; the two velocity
factors in advection and the spatial derivative give
$\nu\nu^{-1/2}$; viscosity and two spatial derivatives give
$\nu\sqrt\nu\nu^{-1}$; and the pressure gradient gives
$\nu\nu^{-1/2}$. The divergence is unchanged. The Jacobian
$dx=\nu^{3/2}dy$ gives
$\|u_\nu\|_2^2=\nu^{5/2}\|u\|_2^2$ and
$\nu\int\|\nabla u_\nu\|_2^2dt
=\nu^{5/2}\int\|\nabla u\|_2^2dt$ on the same time interval.
In these rescaled fields the physical coordinates satisfy
$z=\sqrt\nu Q^D\xi$, $\rho^2=2\nu QX$ and
$\tau=Q(1-\xi^2)$. Thus $\partial_z$ gains $\nu^{-1/2}$
and the radial scalar Laplacian gains $\nu^{-1}$ in
\eqref{eq:ns-profile-operators-exact}. This is the source's specified
viscosity transport, with its time coordinate unchanged.

\subsection{The material base and the complete two-variable conic connection}
Work first over $\mathbb C$. Retain the original coordinates $s,b,c,r$
and $P=cr^3-2r^2+br-4s$. The original material base derivation and
its simple-root lift are
\[
 U_B=\partial_s+6c\partial_b,\qquad
 \mathcal D=\partial_s+6c\partial_b
        +\frac{4-6cr}{3cr^2-4r+b}\partial_r.
\]
The three contributions to $\mathcal DP$ are $-4$, $6cr$, and
$4-6cr$, so this is a derivation of the localization at $P_r$.
Under the real base identification
\begin{equation}\label{eq:ns-original-base-identification}
 s=\tau=1-t,\qquad b=z,\qquad c=c,
\end{equation}
$U_B$ becomes $-\partial_t+6c\partial_z$. This identity concerns
the stated material base; the leading physical material derivative
has instead the explicitly proved profile formula
\eqref{eq:ns-source-material-profile}. On $c=0$ the original
arithmetic operator is exactly $-\partial_s^2/4=-\partial_t^2/4$.
The previously established pullback by the original polynomial $F$
composes with \eqref{eq:ns-original-base-identification}; its two
original Euclidean diffusion coefficients are unchanged.

On $c=0$, retain the complete cover
\begin{gather*}
 \beta=b/4,\qquad \eta=r-\beta,\qquad
 m=b^2/16-2s,\qquad n=m+4,\\
 \mathcal B=\mathbb C[s,b,(mn)^{-1}],\qquad
 \mathcal E=\mathcal B[\eta,p]/(\eta^2-m,p^2-n).
\end{gather*}
The original branch factors are $\delta=16m=b^2-32s$ and
$L=n=4+b^2/16-2s$. The localization removes both and only these
two factors. The ring is free of rank four, with bases
$(1,\eta,p,\eta p)$ and $(1,r,p,rp)$; monic division proves
existence and uniqueness of the four coefficients. There are inverse
maps to the localized original Laurent presentation, given by
\begin{equation}\label{eq:ns-conic-laurent-map}
 q=(p+\eta)/2,\quad q^{-1}=(p-\eta)/2,
 \quad\eta=q-q^{-1},\quad p=q+q^{-1},\quad
 s=b^2/32-\tfrac12(q-q^{-1})^2.
\end{equation}
The two displayed images of $q,q^{-1}$ have product
$(n-m)/4=1$; sums and differences prove both inverse compositions.
Thus the arithmetic $q$ in this formula is never set equal to $Q$.

There are unique commuting derivations $\nabla_s,\nabla_b$ extending
$\partial_s,\partial_b$ to $\mathcal E$. On generators they are
\begin{align*}
 \nabla_s\eta&=-1/\eta,&\nabla_sp&=-1/p,&
 \nabla_sq&=-q/(\eta p),\\
 \nabla_b\eta&=b/(16\eta),&\nabla_bp&=b/(16p),&
 \nabla_bq&=bq/(16\eta p).
\end{align*}
Differentiating the two defining quadratic relations proves these
formulas and their uniqueness, since $2\eta$ and $2p$ are units.
Their commutator is a derivation killing $s,b$; differentiation of
the same quadratic relations then forces it to kill $\eta,p$.
This proves flatness without assuming that arbitrary differentiation
descends to an unlocalized quotient.

Here are the full matrices in the retained basis $(1,r,p,rp)$.
Column vectors are coefficient vectors, and
$\nabla_i=\partial_iI_4+\Gamma_i$:
\begin{align*}
 \Gamma_s&=\begin{pmatrix}
 0&\beta/m&0&0\\0&-1/m&0&0\\
 0&0&-1/n&\beta/m\\0&0&0&-1/m-1/n
 \end{pmatrix},\\
 \Gamma_b&=\begin{pmatrix}
 0&1/4-b\beta/(16m)&0&0\\0&b/(16m)&0&0\\
 0&0&b/(16n)&1/4-b\beta/(16m)\\
 0&0&0&b/(16m)+b/(16n)
 \end{pmatrix}.
\end{align*}
For example $\nabla_br=1/4+b\eta/(16m)$ and
$\nabla_b(rp)=p/4+(b/(16m)+b/(16n))\eta p
+b\beta p/(16n)$, which give the second and fourth columns.
The other columns follow directly from the generator formulas.
The inverse change of coefficient bases is explicit:
\[
 T=\begin{pmatrix}1&\beta&0&0\\0&1&0&0\\
 0&0&1&\beta\\0&0&0&1\end{pmatrix},\qquad
 T^{-1}=\begin{pmatrix}1&-\beta&0&0\\0&1&0&0\\
 0&0&1&-\beta\\0&0&0&1\end{pmatrix}.
\]
In the $(1,\eta,p,\eta p)$ basis write
\[
 H_s=\operatorname{diag}(0,-m^{-1},-n^{-1},-m^{-1}-n^{-1}),
 \qquad H_b=-\tfrac b{16}H_s.
\]
Then $\Gamma_i=T^{-1}H_iT+T^{-1}(\partial_iT)$, which
also verifies that the $b$-dependent basis derivative has been retained.
For clarity the entire second axial derivative in this basis is
\[
 \nabla_b^2=\partial_b^2I_4+2H_b\partial_b+J_b,
\]
where $J_b$ is diagonal with entries
\begin{align*}
 (J_b)_{00}&=0,\\
 (J_b)_{11}&=\frac1{16m}-\frac{b^2}{256m^2},\qquad
 (J_b)_{22}=\frac1{16n}-\frac{b^2}{256n^2},\\
 (J_b)_{33}&=\frac1{16}\left(\frac1m+\frac1n\right)
 -\frac{b^2}{256}\left(\frac1{m^2}+\frac1{n^2}\right)
 +\frac{b^2}{128mn}.
\end{align*}
These entries are $\partial_bH_b+H_b^2$, using
$\partial_bm=\partial_bn=b/8$. Similarly
\[
 \nabla_s^2=\partial_s^2I_4+2H_s\partial_s+J_s,
 \quad J_s=\operatorname{diag}
 \left(0,-m^{-2},-n^{-2},-m^{-2}-n^{-2}+2/(mn)\right).
\]
The formulas follow from $\partial_sm=\partial_sn=-2$.
Returning to the original basis replaces each second-order constant
matrix by $\partial_i\Gamma_i+\Gamma_i^2$ and each first-order
matrix by $2\Gamma_i$, so no basis-dependent second derivative is lost.

Let the physical real base be
\[
 \Omega=\{(\rho,z,t):\rho>0,\ t<1,\ m=z^2/16-2(1-t),\ mn\ne0\}.
\]
Tensoring $\mathcal E$ over its indicated base map with complex smooth
functions on $\Omega$ gives a rank-four complex algebra bundle.
On this bundle physical $\partial_t=-\nabla_s$,
$\partial_z=\nabla_b$, and $\partial_\rho$ differentiates the smooth
coefficients and kills $\eta,p$. Consequently its complete scalar
heat and arithmetic operators are
\begin{equation}\label{eq:ns-flat-heat-connection}
 \mathcal N_\nu=-\nabla_s-\nu\bigl(
 \partial_\rho^2+\rho^{-1}\partial_\rho+\nabla_b^2\bigr),
 \qquad \mathcal H=-\tfrac14\nabla_s^2,
 \qquad [\mathcal N_\nu,\mathcal H]=0.
\end{equation}
The commutator vanishes because the two lifted derivations commute,
all transverse derivatives commute with them, and $\nu$ is constant.
The coefficient inclusion $f\mapsto f\,1$ is an algebra map and
intertwines every displayed physical derivative. Inserting the actual
profiles into it therefore carries \eqref{eq:ns-profile-heat-full},
\eqref{eq:ns-profile-arithmetic-full}, and the three exact residual
identities to this bundle. This is the explicit differential
correspondence; it does not identify the four cover coordinates with
four independently constructed physical velocities.

\subsection{The physical real structure and a finite negative coefficient}
On the open physical region $-4<m<0$, put $a_*=\sqrt{-m}$ and
$p_*=\sqrt{m+4}$ using the positive real square roots.
The four sections of the complex cover are exactly
\[
 \eta=\epsilon i a_*,\quad p=\epsilon' p_*,\quad
 r=b/4+\epsilon i a_*,\quad
 q=\tfrac12(\epsilon'p_*+\epsilon i a_*),\qquad
 \epsilon,\epsilon'\in\{1,-1\}.
\]
Their squared moduli are $(m+4-m)/4=1$, and the maps
\eqref{eq:ns-conic-laurent-map} verify every equation. The four
sections are distinct since $a_*,p_*>0$. Ordinary complex conjugation
fixes $s,b$, fixes $p$, changes $\eta$ to $-\eta$, and sends $q$ to
$q^{-1}$. This is the real structure over the physical real base.
There is no real value of $q$ on this region: a real nonzero $q$ would
give $m=(q-q^{-1})^2\geq0$. The branch boundaries $m=0$ and $m=-4$
remain excluded from the connection, even though limits of the scalar
polynomials below exist there.

On the source's actual growth path $z=0$ and
$\rho=\sqrt{2X_{\rm in}\tau}$, retain $Q=\tau$, $\xi=0$,
$X=X_{\rm in}>0$. The base map gives $b=0,s=\tau$, so
$m=-2\tau$. For $0<\tau<2$ choose the sheet
\begin{equation}\label{eq:ns-unit-circle-growth-map}
 \theta=\arcsin\sqrt{\tau/2},\quad
 q=e^{i\theta},\quad \eta=i\sqrt{2\tau},\quad
 p=\sqrt{4-2\tau},\quad r=i\sqrt{2\tau}.
\end{equation}
Here $\theta\in(0,\pi/2)$ and the formulas agree with the four-sheet
description. In particular the source concentration $Q\downarrow0$
maps to the finite arithmetic point $q\to1$, with the original branch
point retained as a limit.

For the original GCT coefficient, every factor is retained:
\begin{align}
 C(q)&=q^{10}-q^4-q^{-4}+q^{-10}
       =(q-q^{-1})^2[7]_q[3]_q,\notag\\
 C(e^{i\theta})&=-4\sin(7\theta)\sin(3\theta),\notag\\
 C\big|_{b=0,s=\tau}
 &=-2\tau(3-2\tau)(7-28\tau+28\tau^2-8\tau^3)
   =\tau B(\tau),\label{eq:ns-finite-negative-polynomial}\\
 B(\tau)&=-42+196\tau-280\tau^2+160\tau^3-32\tau^4.
       \notag
\end{align}
The trigonometric equality follows by substituting $q=e^{i\theta}$
and applying $2\cos(10\theta)-2\cos(4\theta)
=-4\sin(7\theta)\sin(3\theta)$. For the polynomial equality,
$[3]_q=p^2-1$ and $[7]_q=p^6-5p^4+6p^2-1$ follow by multiplying
the corresponding finite Laurent sums by $q-q^{-1}$. Substituting
$p^2=4-2\tau$ and $(q-q^{-1})^2=-2\tau$ gives the displayed
factored polynomial and its full expansion. Therefore
\[
 C(q)<0\quad\hbox{for}\quad
 0<\tau<2\sin^2(\pi/7),\qquad
 \lim_{\tau\downarrow0}\frac{C(q)}\tau=-42.
\]
Both sines are strictly positive on the stated interval because
$0<\theta<\pi/7$ and $0<3\theta<3\pi/7<\pi$.
This proves a finite negative value and its exact vanishing order.

The precise scalar-germ transport of this coefficient is also explicit.
Fix $0<h<1/100$, $e_0>0$, and any bounded real function
$\varepsilon$ on $(0,\tau_0)$. Define the actual scalar function
\[
 g_\varepsilon(\tau)=\tau^{-A}
          (e_0+\tau^{2h}\varepsilon(\tau)).
\]
Multiplication by the computed GCT coefficient gives the exact identity
\begin{equation}\label{eq:ns-negative-germ-transport}
 C(q)g_\varepsilon(\tau)
 =\tau^D B(\tau)(e_0+\tau^{2h}\varepsilon(\tau)),\qquad
 \lim_{\tau\downarrow0}\tau^{-D}C(q)g_\varepsilon(\tau)=-42e_0.
\end{equation}
This follows from $1-A=D$ and \eqref{eq:ns-finite-negative-polynomial}.
Since $B(\tau)+42$ is divisible by $\tau$, and $0<2h<1$, the
right side equals $-42e_0\tau^D+O(\tau^{D+2h})$.
It is negative for sufficiently small $\tau$ and tends to zero,
since $D>0$. Multiplication raises the displayed leading exponent by
exactly one. In the formal power-series ring $\mathbb R[[\tau]]$,
the same exact polynomial has order one, because $B(0)=-42\ne0$.
No identification of $\tau$ with a prime number or of this order with
a prime-adic valuation is made.

The source's equation (10.21) reports its completed velocity component
as $\tau^{-A}(e_0+O(\tau^{2h}))$ on this path. The exact map
\eqref{eq:ns-negative-germ-transport} acts on the full stated class of
such scalar germs; identifying a member with that completed solution
retains the independent-verification status of the source theorem.
The negative coefficient result above uses only the proved conic map
and polynomial identities. It gives no asymptotic circuit-size bound,
plethysm multiplicity obstruction, or nonstandard-RH eigenvalue claim.

Finally, the source pressure normalization supplies a second exact
sign statement with a different specified observable:
$\Pi(X,\xi)=-\int_X^\infty E(y,\xi)^2/(2y)\,dy\leq0$
whenever that integral is finite. It is strictly negative if the
nonnegative integrand is positive on a set of positive measure.
Replacing $p$ by $p+k(t)$ leaves its spatial gradient and the momentum
equation unchanged, but need not preserve a chosen compact-support or
infinity normalization. Thus the displayed pressure sign belongs to
that normalization. The coefficient sign
\eqref{eq:ns-finite-negative-polynomial} has instead been proved for
the fixed original Laurent polynomial on the stated unit-circle sheet.
